\section{Induction}
\label{sec:induction-principles}

\begin{topicbox}{Strong Induction}
\begin{exposition}
Ordinary induction propagates a property from each element to its immediate
successor. Strong induction allows the inductive step to use the property at
every predecessor, not just the immediate one. This broader hypothesis makes
strong induction strictly more convenient for arguments in which the successor
step requires information from elements earlier than the immediate predecessor.
\end{exposition}

\begin{tcolorbox}[colback=thmbox, colframe=thmborder, arc=2pt,
  left=6pt, right=6pt, top=4pt, bottom=4pt,
  title={\small\textbf{Theorem (Strong Induction on a Peano System)}},
  fonttitle=\small\bfseries]
\begin{theorem}[Strong Induction on a Peano System]
\label{thm:strong-induction-on-peano-system}
Let $(P,S,1)$ be a Peano system, and let $\varphi$ be a predicate on $P$.
If
\[
\forall n \in P\;
\bigl(\forall k \in P\;(k < n \Longrightarrow \varphi(k))\bigr)
\Longrightarrow \varphi(n),
\]
then
\[
\forall n \in P\;\varphi(n).
\]
\hyperref[prf:strong-induction-on-peano-system]{Go to proof.}
\end{theorem}
\end{tcolorbox}

\begin{remark*}[Standard quantified statement]
\[
\begin{aligned}
&\text{Let } (P,S,1) \text{ be a Peano system and let } \varphi \text{ be a predicate on } P.\\
&\Bigl(\forall n \in P\;
\bigl(\forall k \in P\;(k < n \Longrightarrow \varphi(k))\bigr)
\Longrightarrow \varphi(n)\Bigr)
\Longrightarrow
\forall n \in P\;\varphi(n).
\end{aligned}
\]
\end{remark*}

\begin{remark*}[Contrapositive quantified statement]
\[
\begin{aligned}
&\text{Let } (P,S,1) \text{ be a Peano system and let } \varphi \text{ be a predicate on } P.\\
&\exists n \in P\;\neg\varphi(n)
\Longrightarrow
\exists n \in P\;
\Bigl(\bigl(\forall k \in P\;(k < n \Longrightarrow \varphi(k))\bigr)
\land \neg\varphi(n)\Bigr).
\end{aligned}
\]
\end{remark*}

\begin{remark*}[Interpretation]
Strong induction grants the inductive step access to the entire history below
$n$, not just the single predecessor. The contrapositive gives the least
counterexample principle: if any element fails $\varphi$, there is a least
one, and that least one has $\varphi$ holding at all its predecessors, which
contradicts the strong-induction hypothesis. This makes strong induction and
least-counterexample arguments two descriptions of the same logical structure.
\end{remark*}

\begin{remark*}[Dependencies]
\begin{itemize}
  \item \hyperref[thm:induction-principle-for-peano-system]{Induction Principle for a Peano System}
  \item \hyperref[def:strict-less-than-on-peano-system]{Strict Less Than on a Peano System}
  \item \hyperref[thm:trichotomy-on-peano-system]{Trichotomy on a Peano System}
\end{itemize}
\end{remark*}
\end{topicbox}

\begin{topicbox}{Induction from an Arbitrary Base}
\begin{exposition}
Ordinary and strong induction each begin at the distinguished first element
$1$. Many arguments require starting at a later element. Induction from an
arbitrary base $m$ restricts the claim to elements at or above $m$ and
starts the inductive process there.
\end{exposition}

\begin{tcolorbox}[colback=thmbox, colframe=thmborder, arc=2pt,
  left=6pt, right=6pt, top=4pt, bottom=4pt,
  title={\small\textbf{Theorem (Induction from an Arbitrary Base)}},
  fonttitle=\small\bfseries]
\begin{theorem}[Induction from an Arbitrary Base]
\label{thm:induction-from-arbitrary-base}
Let $(P,S,1)$ be a Peano system, let $m \in P$, and let $\varphi$ be a
predicate on $P$. If
\[
\varphi(m)
\qquad\text{and}\qquad
\forall n \in P\;\bigl(n \ge m \land \varphi(n) \Longrightarrow \varphi(S(n))\bigr),
\]
then
\[
\forall n \in P\;\bigl(n \ge m \Longrightarrow \varphi(n)\bigr).
\]
\hyperref[prf:induction-from-arbitrary-base]{Go to proof.}
\end{theorem}
\end{tcolorbox}

\begin{remark*}[Standard quantified statement]
\[
\begin{aligned}
&\text{Let } (P,S,1) \text{ be a Peano system, } m \in P,
\text{ and } \varphi \text{ a predicate on } P.\\
&\Bigl(\varphi(m) \land
\forall n \in P\;\bigl(n \ge m \land \varphi(n)
\Longrightarrow \varphi(S(n))\bigr)\Bigr)
\Longrightarrow
\forall n \in P\;\bigl(n \ge m \Longrightarrow \varphi(n)\bigr).
\end{aligned}
\]
\end{remark*}

\begin{remark*}[Interpretation]
Apply ordinary induction to the shifted predicate $\psi(n) \coloneqq
(n \ge m \Longrightarrow \varphi(n))$, which holds vacuously below $m$
and reduces to $\varphi$ above $m$. This theorem is the formal license for
all arguments of the form ``for all $n \ge m$, proceed by induction.''
\end{remark*}

\begin{remark*}[Dependencies]
\begin{itemize}
  \item \hyperref[thm:induction-principle-for-peano-system]{Induction Principle for a Peano System}
  \item \hyperref[def:less-than-or-equal-on-peano-system]{Less Than or Equal on a Peano System}
\end{itemize}
\end{remark*}
\end{topicbox}
